\section{Timing Analysis with the Logic Analyzer}
\label{sec:timing}

\begin{frame}{Measuring timing in Wokwi}
  \begin{itemize}
    \item The patched FreeRTOS library (\texttt{lib/FreeRTOS}) is used together with the Wokwi logic analyzer.
    \item Digital pins carry debug signals: one pin per task, one per ISR.
    \item The recorded signals are opened in a web-based waveform viewer.
    \item Reference project: \url{https://wokwi.com/projects/446797093390429185}
  \end{itemize}
\end{frame}

\scriptonly{%
The reference project (interrupt-triggered semaphore) uses pin~2 as button input (INT0), pin~13 and pin~12 for the two LEDs, and pins~10 and~11 as debug outputs. The logic analyzer is connected to D0--D7.
}

\begin{frame}[fragile]{The patched library: scheduler-driven trace pins}
  \begin{itemize}
    \item The trace hooks \texttt{traceTASK\_SWITCHED\_IN()} and \texttt{traceTASK\_SWITCHED\_OUT()} are defined in \texttt{Arduino\_FreeRTOS.h}.
    \item Each task carries a \emph{application tag}: the pin number, set with \texttt{vTaskSetApplicationTaskTag()}.
    \item Scheduler switches task \emph{in} (Running) $\Rightarrow$ pin HIGH. Task switched \emph{out} (Ready or Blocked) $\Rightarrow$ pin LOW.
  \end{itemize}
\begin{lstlisting}
#define traceTASK_SWITCHED_IN()  \
    analogWrite((int)pxCurrentTCB->pxTaskTag, 255)
#define traceTASK_SWITCHED_OUT() \
    analogWrite((int)pxCurrentTCB->pxTaskTag, 0)

pinMode(10, OUTPUT);
vTaskSetApplicationTaskTag(handle, (TaskHookFunction_t)10);
\end{lstlisting}
\end{frame}

\scriptonly{%
The pin therefore shows exactly the time a task is \emph{Running}, as decided by the scheduler, and not merely the time between two statements in the task code. Preemption is visible as a gap in the pulse of the lower-priority task while the pulse of the higher-priority task is high. The application tag must be set before the scheduler starts or at least before the task is first switched in. Two things to keep in mind: tasks without a tag have the tag \texttt{NULL}, so the macros drive pin~0 for them (this includes the idle task, and pin~0 is the serial RX pin on the Uno); and the pin has to be configured with \texttt{pinMode()} as output.
}

\begin{frame}{What to read from the waveform}
  \begin{itemize}
    \item \textbf{Latency} $T_L$: event edge $\rightarrow$ start of task pin.
    \item \textbf{Execution time} $T_V$: width of the task pulse.
    \item \textbf{Reaction time} $T_R$: event edge $\rightarrow$ end of task pulse.
    \item \textbf{Jitter}: variation of period and reaction time over many samples.
    \item \textbf{Preemption}: the lower-priority pulse is interrupted by a higher-priority one.
  \end{itemize}
\end{frame}
